\section{Introduction}
\label{sec:intro}

Particle imaging detectors infer the properties of fundamental particles from spatially resolved measurements of their interactions with matter. The resulting images can differ radically between experiments: liquid argon time projection chambers (LArTPCs) directly image three-dimensional ionization trajectories~\citep{Abi2020}; collider TPCs such as sPHENIX observe dense, curved trajectories produced in a magnetic field~\citep{klest2022compacttpc}; and water Cherenkov detectors such as Super-Kamiokande (SK) instead record patterns of Cherenkov photons projected onto a detector surface~\citep{fukuda2003superkamiokande}. Despite originating from related particle processes, these modalities differ substantially in geometry, density, sensing mechanism, and characteristic physical structure (Fig.~\ref{fig:overview}).

Modern reconstruction methods consequently tend to be highly detector-specific, combining specialized representations, architectural inductive biases, and task-specific objectives~\citep{Acciarri2018,Abratenko_2022,drielsma2021scalableendtoenddeeplearningbaseddata,park2025fm4npp}. This contrasts with the foundation model paradigm, where representations learned from abundant unlabeled data can be reused across many downstream tasks with little supervision~\citep{oquab2024dinov2learningrobustvisual,wu2025sonataselfsupervisedlearningreliable}. \cite{panda} recently demonstrated this principle at the sensor level for LArTPC data, but whether the same representation-learning mechanism can operate across fundamentally different particle detectors remains unknown.

Here we introduce \textbf{Panda V2}, an updated backbone architecture, and independently pre-train the same architecture and self-distillation objective on LArTPC, sPHENIX TPC, and SK-like water Cherenkov data. Only hyperparameters tied directly to the numerical and spatial scale of each input are changed between modalities. The objective contains no detector-specific propagation or reconstruction rules, and we apply strong generic spatial augmentations (rotations, translations, jittering). Nevertheless, the learned representations develop clear semantic organization across all three modalities (Fig.~\ref{fig:overview} c-e).

We evaluate the learned representations by asking both {what reconstruction they enable} and {what physical structure they encode}. Using only \textbf{1,000 labeled events} per downstream task, lightweight readouts recover point-, particle-, and event-level observables across all three detectors, while remaining competitive with specialized reconstruction methods trained with substantially more supervision. We additionally find simple latent directions associated with the ``arrow of time" for individual particles in LArTPCs and transverse momentum in sPHENIX. These results suggest that the machinery required for foundation-model pre-training on sensor-level particle reconstruction can be general, and point toward unified systems that can ingest heterogeneous data from one or more experiments.